\documentclass[11pt]{article}
\usepackage[T5]{fontenc}
\usepackage[utf8]{inputenc}
\usepackage[vietnamese,english]{babel}
\usepackage[a4paper,top=2cm,bottom=2cm,left=2.7cm,right=2.7cm]{geometry}
\usepackage{amsmath}
\usepackage{amssymb}      % \square, \boxtimes cho tickbox kế hoạch
\usepackage{graphicx}
\usepackage{booktabs}
\usepackage{array}
\newcolumntype{P}[1]{>{\centering\arraybackslash}p{#1}}  % cột căn giữa, rộng cố định
\usepackage[table]{xcolor}
\usepackage{caption}
\usepackage{float}
\usepackage[square,numbers]{natbib}
\usepackage[colorinlistoftodos]{todonotes}
\usepackage{listings}
\usepackage{tikz}
\usetikzlibrary{arrows.meta, positioning, calc, fit, backgrounds}
\usepackage{hyperref}
\definecolor{lightblue}{RGB}{173, 216, 230}

% Ánh xạ vài ký tự Unicode đặc biệt để biên dịch được bằng pdflatex (T5+inputenc):
% 【 】 = ngoặc trích dẫn của hệ thống; �� = nhãn nguồn web.
\DeclareUnicodeCharacter{3010}{[}
\DeclareUnicodeCharacter{3011}{]}
\DeclareUnicodeCharacter{1F310}{(web)}

\hypersetup{
    colorlinks=true,
    linkcolor=blue,
    filecolor=magenta,
    urlcolor=cyan,
    citecolor=blue,
}

% Tickbox cho phần kế hoạch: \cbDone = đã xong, \cbTodo = chưa/đang làm.
\newcommand{\cbDone}{\item[$\boxtimes$]}
\newcommand{\cbTodo}{\item[$\square$]}

% Style cho các sơ đồ luồng (TikZ) trong phần Method.
\definecolor{boxblue}{RGB}{225,238,252}
\definecolor{boxgreen}{RGB}{223,244,228}
\definecolor{boxamber}{RGB}{253,243,222}
\definecolor{boxgray}{RGB}{238,238,242}
\tikzset{
  fbox/.style={rectangle, rounded corners=2pt, draw=black!55, fill=boxblue,
               align=center, font=\footnotesize, minimum height=0.95cm,
               inner sep=4pt, text width=2.5cm},
  dbox/.style={fbox, fill=boxgreen},
  abox/.style={fbox, fill=boxamber},
  gbox/.style={fbox, fill=boxgray},
  farr/.style={-{Stealth[length=2.2mm]}, semithick, draw=black!65},
  flab/.style={font=\scriptsize\itshape, fill=white, inner sep=1pt},
}

\setlength{\parskip}{2pt}
\renewcommand{\baselinestretch}{1.05}

\title{\textbf{Assistant hỏi đáp thủ tục hành chính công Việt Nam ứng dụng Agentic RAG}}
\author{
    Nguyễn Xuân Phúc\textsuperscript{1},
    Nguyễn Minh Triết\textsuperscript{1},
    Nguyễn Ngọc Quí\textsuperscript{1,*}
    Đỗ Trọng Hợp\textsuperscript{1,*},\\
    \vspace{0.3cm}
    \small \textsuperscript{1}\textit{Trường Đại học Công nghệ Thông tin, ĐHQG-HCM, TP. Hồ Chí Minh}
}
\date{}
\begin{document}

\maketitle
\let\thefootnote\relax\footnotetext{\textsuperscript{*}Giảng viên hướng dẫn}

\begin{abstract}
Thủ tục hành chính công ở Việt Nam bao gồm hàng nghìn quy trình phức tạp, có nhiều thành phần hồ sơ và thường xuyên thay đổi, gây khó khăn cho người dân và doanh nghiệp trong quá trình tra cứu. Báo cáo này trình bày việc xây dựng một trợ lý hỏi-đáp tiếng Việt chuyên biệt cho miền thủ tục hành chính công, dựa trên kiến trúc Agentic RAG (Retrieval-Augmented Generation). Kho tri thức của hệ thống được thu thập và chuẩn hóa trực tiếp từ Cổng Dịch vụ công Quốc gia, lưu trữ dưới dạng cơ sở dữ liệu vector. Lõi suy luận của hệ thống sử dụng cơ chế tác tử thông minh (agent) có khả năng tự động tra cứu dữ liệu nội bộ và kết hợp tìm kiếm web khi cần thiết. Từ đó, trợ lý có thể tổng hợp câu trả lời chính xác, kèm theo tính năng trích dẫn nguồn rõ ràng. Ngoài ra, ứng dụng còn được tích hợp các chức năng nâng cao như trực quan hóa quá trình suy luận, gợi ý câu hỏi tiếp theo và sinh sơ đồ tư duy cho từng thủ tục trên một giao diện web thân thiện. Báo cáo tập trung mô tả động lực nghiên cứu, các công trình liên quan và kế hoạch triển khai hệ thống, tạo tiền đề cho các bước hoàn thiện và đánh giá chuyên sâu tiếp theo.
\end{abstract}

% ===========================================================================
\section{Introduction}
% ===========================================================================

\subsection{Bối cảnh và động lực}

Cải cách thủ tục hành chính là một trong những trụ cột của tiến trình xây dựng Chính phủ số ở Việt Nam. Tuy nhiên, hệ thống thủ tục hành chính công hiện hành gồm hàng nghìn thủ tục thuộc nhiều lĩnh vực và nhiều cấp thực hiện (xã/phường, tỉnh/thành phố, bộ/trung ương); mỗi thủ tục lại bao gồm thành phần hồ sơ, trình tự thực hiện, cách thức nộp, lệ phí, thời hạn giải quyết và căn cứ pháp lý riêng. Các thông tin này phân tán, được diễn đạt bằng văn phong pháp lý, và thường xuyên thay đổi theo các thông tư, nghị định mới. Hệ quả là người dân và doanh nghiệp dễ nộp sai cơ quan, thiếu giấy tờ, phải đi lại nhiều lần gây lãng phí thời gian và chi phí xã hội.

Cổng Dịch vụ công Quốc gia~\citep{dichvucong} tập trung dữ liệu thủ tục, nhưng việc tra cứu chủ yếu dựa trên từ khoá: người dùng phải biết gần đúng tên thủ tục, và hệ thống không hỗ trợ hỏi đáp bằng ngôn ngữ tự nhiên, không tổng hợp hay đối chiếu thông tin giữa các thủ tục liên quan. Các mô hình ngôn ngữ lớn (LLM) như GPT-4~\citep{openai2023gpt4} giải quyết tốt phần giao tiếp ngôn ngữ tự nhiên, nhưng khi hỏi về dữ kiện hành chính cụ thể (mã mẫu đơn, số văn bản, mức lệ phí, thành phần hồ sơ) chúng dễ ảo giác sinh ra thông tin nghe hợp lý nhưng sai~\citep{ji2023hallucination}. Trong một miền đòi hỏi độ chính xác cao và phải truy được nguồn như thủ tục hành chính, rủi ro này khiến sai lệch trong các thủ tục hành chính.

Retrieval-Augmented Generation (RAG) nối LLM với một kho tri thức ngoài, giúp giảm ảo giác và cập nhật thông tin mà không cần huấn luyện lại~\citep{lewis2020rag,gao2023ragsurvey}. Tuy vậy, RAG tĩnh truy hồi một lần rồi sinh câu trả lời bộc lộ hạn chế trong bài toán của nhóm: một truy vấn diễn đạt kém sẽ làm sai lệch ngữ cảnh vốn là thứ cần thiết nhất để trả lời câu hỏi; hệ thống không tự đánh giá mức độ liên quan để sửa sai; và nhiều thông tin bám sát với thực tế mà người dân thực sự cần (địa chỉ cụ thể của cơ quan tiếp nhận, giờ làm việc, đường dây nóng theo địa bàn) vốn không nằm trong văn bản mô tả thủ tục. Cách tiếp cận Agentic RAG mà tiêu biểu là cách tiếp cận suy luận: ReAct~\citep{yao2023react}, kết hợp suy luận và hành động cho phép xem LLM như một agent biết tự quyết định: chủ động truy hồi lại tập corpus, tự đánh giá thông tin đã đủ để suy luận hay chưa, và bổ sung tìm kiếm thêm thông tin trên web khi cần theo cơ chế tự sửa của Corrective-RAG~\citep{yan2024crag, asai2024selfrag}. Cuối cùng, để người dùng tin tưởng vào truy vấn và kiểm chứng được, trợ lý  agent phải trích dẫn nguồn ngay trong câu trả lời, tô sáng đúng đoạn bằng chứng, và minh bạch hoá chuỗi suy luận của agent.

\subsection{Phát biểu bài toán}

Với bối cảnh như vậy, nhómđặt ra mục tiêu của đồ án là xây dựng một trợ lý hỏi đáp cho thủ tục hành chính công của Việt Nam, thoả các yêu cầu tất yếu: \textbf{(i)} có căn cứ các khẳng định phải dựa trên kho 5.208 thủ tục
thu thập từ Cổng Dịch vụ công Quốc gia, hoặc nguồn web được đánh dấu rõ là đã thẩm định; \textbf{(ii)} agentic tự tra cứu corpus trước, bổ sung web khi corpus thiếu, hỗ trợ hội
thoại dài bằng việc cung cấp bộ nhớ ngữ cảnh; \textbf{(iii)} minh bạch trích dẫn nội dòng kèm tô sáng bằng chứng và hiển thị quá trình suy luận của agent; và \textbf{(iv)} linh hoạt triển khai chạy được với nhiều nhà cung cấp LLM khác nhau (Azure, OpenAI, mô hình mã nguồn mở phục vụ qua vLLM).

\subsection{Đóng góp}

Báo cáo có các đóng góp chính sau:
\begin{itemize}
  \item \textbf{Pipeline dữ liệu end-to-end có khả năng cập nhật: } Xây dựng một crawler có khả năng dịch ngược (reverse-engineer) các API ẩn từ Cổng Dịch vụ công Quốc gia để thu thập trọn bộ 5.208 thủ tục. Dữ liệu thô sau đó được xử lý qua parser để chuẩn hoá từ định dạng JSON sang dạng corpus Markdown có cấu trúc (gồm frontmatter metadata và sáu mục: trình tự, cách thức, thành phần hồ sơ, điều kiện, căn cứ pháp lý, kết quả), trước khi được nhúng và nạp (ingest) vào cơ sở dữ liệu lai (hybrid vector/document store).
  \item \textbf{Kiến trúc Agentic RAG hai giai đoạn chuyên biệt: } Giai đoạn một triển khai vòng lặp ReAct với hai công cụ: \texttt{docsearch} (tra cứu corpus, được ưu tiên và tích hợp cơ chế lọc độ liên quan) và \texttt{web\_search} (công cụ dự phòng qua Brave Search, thiết kế theo mô hình Corrective-RAG). Giai đoạn hai chịu trách nhiệm tổng hợp câu trả lời và tạo trích dẫn nội dòng (inline citation).
  \item \textbf{Cung cấp dẫn chứng uy tín
  : } Trực quan hoá chuỗi Thought/Action của agent theo thời gian thực; cung cấp trích dẫn nội dòng \texttt{(Ví dụ:【n】)} kết hợp highlight chính xác đoạn bằng chứng gốc; dán nhãn phân biệt các nguồn web chưa qua kiểm chứng. Hệ thống cũng tích hợp khả năng sinh sơ đồ tư duy cho từng thủ tục, gợi ý câu hỏi tiếp theo và tự động viết lại truy vấn (query rewriting) bám sát ngữ cảnh hội thoại.
\end{itemize}

Phần còn lại của báo cáo được cấu trúc như sau: Mục~\ref{sec:related} tổng quan các công trình nghiên cứu liên quan. Mục~\ref{sec:method} trình bày chi tiết phương pháp đề xuất. Cuối cùng, Mục~\ref{sec:plan} mô tả kế hoạch triển khai trong bốn tuần cùng tiến độ thực hiện.

% ===========================================================================
\section{Related Work}
\label{sec:related}
% ===========================================================================

\paragraph{Nền tảng RAG và hỏi--đáp miền mở.}
RAG~\citep{lewis2020rag} kết hợp một bộ truy hồi với một bộ sinh để trả lời các tác vụ đòi
hỏi nhiều tri thức, và đã được tổng hợp lại trong khảo sát của \citet{gao2023ragsurvey}
theo ba thế hệ (naive, advanced, modular). Các kiến trúc liên quan như Fusion-in-Decoder
\citep{izacard2021fid} cho thấy lợi ích của việc đưa nhiều đoạn truy hồi vào bộ sinh. Hệ
thống của nhóm kế thừa khung này nhưng đặt trong một vòng lặp tác tử thay vì truy hồi
một lần.

\paragraph{Truy hồi: dày đặc, thưa, và xấp xỉ láng giềng.}
Dense Passage Retrieval~\citep{karpukhin2020dpr} và Sentence-BERT~\citep{reimers2019sbert}
mở đường cho truy hồi theo ngữ nghĩa bằng vector nhúng, trong khi BM25~\citep{robertson2009bm25}
vẫn mạnh ở so khớp từ khoá   đặc biệt quan trọng với các thực thể "cứng" như mã mẫu đơn
(CC01) hay số văn bản (Thông tư 17/2024/TT-BCA). Vì vậy nhóm dùng truy hồi lai
ghép (vector + toàn văn). Ở quy mô hàng nghìn tài liệu, tìm kiếm láng giềng gần đúng bằng
đồ thị HNSW~\citep{malkov2020hnsw} cho độ trễ thấp mà vẫn giữ chất lượng.

\paragraph{Suy luận và tác tử dùng công cụ.}
Chain-of-Thought~\citep{wei2022cot} cho thấy việc "suy luận thành lời" cải thiện khả năng
giải quyết vấn đề của LLM. ReAct~\citep{yao2023react} đan xen suy luận với hành
động (gọi công cụ) và quan sát kết quả, tạo nên khung tác tử mà nhóm áp dụng.
Toolformer~\citep{schick2023toolformer} chứng minh LLM có thể tự học cách dùng công cụ, còn
ReWOO~\citep{xu2023rewoo} tách lập kế hoạch khỏi quan sát để giảm số lần gọi LLM   một
phương án thay thế mà kotaemon cũng hỗ trợ và nhóm có thể đối chiếu.

\paragraph{RAG thích nghi và tự sửa.}
Self-RAG~\citep{asai2024selfrag} huấn luyện mô hình tự quyết định khi nào cần truy hồi và tự
phê bình câu trả lời, trong khi Corrective-RAG~\citep{yan2024crag} đánh giá chất lượng tài
liệu truy hồi và kích hoạt tìm kiếm web khi tài liệu nội bộ không đủ tốt. Thiết kế "corpus
trước, web dự phòng" cùng cơ chế cổng lọc độ liên quan của nhóm lấy cảm hứng trực tiếp
từ hai hướng này, nhưng hiện thực bằng prompt trên tác tử ReAct thay vì huấn luyện
lại mô hình.

\paragraph{Tính trung thực, ảo giác và trích dẫn.}
Ảo giác là rào cản cốt lõi khi đưa LLM vào miền pháp lý--hành chính~\citep{ji2023hallucination};
ngoài ra LLM còn "bỏ quên thông tin ở giữa" khi ngữ cảnh dài~\citep{liu2024lostmiddle}. Đây
là lý do nhóm nhấn mạnh trích dẫn nội dòng kèm tô sáng đúng đoạn bằng chứng, đồng thời
ràng buộc bằng prompt để hệ thống chỉ trả lời dựa trên ngữ cảnh đã truy hồi và nói rõ khi
thông tin thiếu.

\paragraph{Khung công cụ và bối cảnh tiếng Việt.}
Đã có những khung RAG mã nguồn mở như kotaemon~\citep{kotaemon2024} giúp rút ngắn thời gian
dựng hệ thống, song phần lớn được thiết kế và đánh giá cho tiếng Anh. So với hỏi--đáp pháp
lý/hành chính tiếng Anh, tài nguyên và hệ thống cho thủ tục hành chính tiếng Việt còn
rất ít. Đóng góp của nhóm nằm ở toàn bộ chuỗi xử lý cho miền này   từ thu thập và
chuẩn hoá dữ liệu, thiết kế prompt theo đặc thù pháp lý hành chính, tới chiến lược suy luận
agentic có trích dẫn   được trình bày ở Mục~\ref{sec:method}.

% ===========================================================================
\section{Methodology}
\label{sec:method}
% ===========================================================================

Hệ thống được tổ chức thành ba thành phần nối tiếp nhau. Thứ nhất, một pipeline xử lý tri thức giúp chuyển đổi dữ liệu thủ tục hành chính thô thành cơ sở dữ liệu có khả năng Retrive (Mục~\ref{sub:data}). Thứ hai, lõi truy hồi và sinh văn bản có trích dẫn (retrieval-generation core) nhận đầu vào là cặp (câu hỏi, kho tri thức) để tạo ra câu trả lời kèm nguồn tham chiếu (Mục~\ref{sub:rag}). Thứ ba, một bộ máy suy luận đa chiến lược (multi-strategy reasoning engine) cho phép linh hoạt thay đổi cách hệ thống lập kế hoạch và gọi công cụ mà không cần cấu hình lại lõi nền tảng (Mục~\ref{sub:reasoning}).

               
\subsection{Thu thập dữ liệu và biểu diễn tri thức}
\label{sub:data}

\paragraph{Nguồn dữ liệu và giao diện truy cập: }
Nguồn tri thức của hệ thống được thu thập hoàn toàn từ Cổng Dịch vụ công Quốc gia\footnote{\url{https://dichvucong.gov.vn}}~\citep{dichvucong}, cổng thông tin chính thức công bố các thủ tục hành chính trên quy mô toàn quốc. Cổng này được xây dựng theo kiến trúc ứng dụng trang đơn (Single-Page Application - SPA), phân tách hoàn toàn giữa giao diện và dữ liệu. Cụ thể, nội dung chi tiết của các thủ tục không được nhúng sẵn trong mã HTML tĩnh mà được nạp động (dynamic loading) thông qua các REST API nội bộ không có tài liệu công khai. Do đó, các kỹ thuật cào dữ liệu truyền thống dựa trên phân tích cây cú pháp HTML (HTML scraping) hầu như không phát huy hiệu quả. Để khắc phục, nhóm tiến hành phân tích lưu lượng mạng (network traffic analysis) của ứng dụng nhằm dịch ngược (reverse-engineer) đặc tả của các endpoint cần thiết. Quá trình này đã xác định được một bộ giao diện ổn định gồm hai endpoint công khai (mang hậu tố \texttt{by-citizen}, thuộc route \texttt{/api/v1}). Các API này không yêu cầu xác thực (authentication) và trả về dữ liệu định dạng JSON theo một cấu trúc chuẩn hóa (gồm mã trạng thái và trường dữ liệu). Bảng~\ref{tab:api} tóm tắt chi tiết hai endpoint này.

\begin{table}[H]
\centering
\small
\renewcommand{\arraystretch}{1.3}
\begin{tabular}{@{}l l p{6.0cm}@{}}
\toprule
Endpoint (POST) & Tham số chính & Vai trò \\
\midrule
\texttt{.../list-all-formality-by-citizen} & \texttt{lastId}, \texttt{limit} & Liệt kê danh mục thủ tục theo phân trang con trỏ \\
\texttt{.../get-formality-by-citizen} & \texttt{id} (UUID) & Tải bản ghi chi tiết đầy đủ của một thủ tục \\
\bottomrule
\end{tabular}
\caption{Hai endpoint REST công khai dùng để thu thập dữ liệu, được khôi phục thông qua phân tích lưu lượng từ Cổng Dịch vụ công Quốc gia.}
\label{tab:api}
\end{table}

\paragraph{Quy trình thu thập dữ liệu: }
Việc thu thập dữ liệu được thực hiện tuần tự qua hai giai đoạn. Giai đoạn duyệt danh mục (Listing) sẽ quét toàn bộ danh sách thủ tục thông qua endpoint đầu tiên bằng cơ chế phân trang con trỏ. Thay vì truy vấn theo chỉ số trang (page index), mỗi request sẽ gửi kèm định danh của phần tử cuối cùng ở trang trước (\texttt{lastId}) làm mốc để máy chủ trả về lô dữ liệu (batch) tiếp theo. Quá trình lặp kết thúc khi máy chủ trả về danh sách rỗng. So với phân trang truyền thống, cơ chế này giúp giảm thiểu rủi ro trùng lặp hoặc bỏ sót dữ liệu trong trường hợp danh mục có sự cập nhật ngay tại thời điểm thu thập. Kết quả của giai đoạn này là một bộ chỉ mục rút gọn (gồm định danh, mã thủ tục, tên, lĩnh vực, cơ quan) cho từng thủ tục. Ở giai đoạn thu thập chi tiết (Detailing), hệ thống lần lượt gọi endpoint thứ hai dựa trên tập định danh để tải bản ghi chi tiết, lưu trữ dưới dạng một tệp JSON độc lập cho mỗi thủ tục.

\paragraph{Đảm bảo tính bền vững và chống bị web ban (Fault Tolerance \& Rate Limiting): } 
Với quy mô truy vấn lên tới hàng nghìn request, quá trình crawler được thiết kế để chịu lỗi tốt và tránh gây quá tải cho máy chủ nguồn. Giữa các request, hệ thống thiết lập một khoảng trễ cơ sở kết hợp nhiễu ngẫu nhiên (random jitter) nhằm dàn đều lưu lượng. Trong trường hợp xảy ra lỗi mạng hoặc lỗi từ phía máy chủ, hệ thống áp dụng cơ chế thử lại với thời gian chờ tăng luỹ thừa (exponential backoff). Nhằm đảm bảo tính toàn vẹn dữ liệu, các bản ghi được ghi xuống ổ đĩa theo cơ chế atomic write xuất ra tệp tạm thời trước khi đổi tên chính thức giúp ngăn chặn việc sinh ra các tệp JSON bị hỏng (corrupted) nếu tiến trình bị ngắt đột ngột. Hệ thống cũng duy trì bộ nhớ trạng thái (state checkpointing) của quá trình phân trang và tập định danh đã tải, cho phép quá trình thu thập có thể tạm dừng và tiếp tục (resume) mà không phải xử lý lại từ đầu. Các định danh thất bại được ghi vào nhật ký (log) riêng để thử lại sau. Toàn bộ dữ liệu thu thập thuộc phạm vi thông tin công khai, phục vụ mục đích nghiên cứu học thuật.

\paragraph{Kết quả thu thập và đặc tả cấu trúc bản ghi: }
Quy trình trên đã thu thập thành công 5.208 bản ghi thủ tục (số liệu tính đến tháng 06/2026). Mỗi bản ghi là một tệp JSON có cấu trúc phức tạp, chứa đầy đủ các thuộc tính của một thủ tục hành chính như: trình tự thực hiện, cách thức (bao gồm thời hạn giải quyết, mức phí và lệ phí), thành phần hồ sơ, yêu cầu--điều kiện, căn cứ pháp lý, kết quả, lĩnh vực, cơ quan thực hiện và đối tượng áp dụng. Điểm đáng lưu ý là trường dữ liệu thành phần hồ sơ, thông tin cốt lõi đối với người dùng cuối  không được đặt ở cấp độ gốc (root level) của JSON mà bị lồng ghép (nested) bên trong mảng các trường hợp thực hiện (\texttt{executionCases}). Việc phân tích chính xác vị trí lồng nhau này là điều kiện tiên quyết để bước chuẩn hoá có thể bóc tách đầy đủ danh mục giấy tờ mà không bị thiếu sót. Khối lượng siêu dữ liệu (metadata) phong phú này cũng cung cấp cơ sở cho các thao tác lọc (filtering) theo lĩnh vực hoặc cấp thẩm quyền, đồng thời hỗ trợ cơ chế truy vết nguồn (source attribution) trong các giai đoạn sau.

\paragraph{Parsing \& Chunking: }
Các bản ghi JSON thô được chuẩn hóa thành định dạng Markdown độc lập (self-contained document), bao gồm hai phần chính: một khối metadata frontmatter (mã thủ tục, lĩnh vực, cấp thực hiện, cơ quan, đối tượng) phục vụ việc lọc và truy vết, cùng phần thân chứa sáu hạng mục nội dung đã qua làm sạch. Thay vì áp dụng kỹ thuật chunking văn bản theo độ dài cố định, nhóm đề xuất phương pháp phân đoạn theo cấu trúc ngữ nghĩa (section-based chunking). Theo đó, mỗi hạng mục của thủ tục sẽ cấu thành một chunk có ngữ cảnh trọn vẹn, được gắn thêm tiền tố định danh \texttt{[Tên thủ tục - Tên mục]} ở đầu. Cách tiếp cận này loại bỏ rủi ro các dữ kiện quan trọng (như thành phần hồ sơ, lệ phí) bị phân mảnh giữa nhiều chunk, đồng thời hỗ trợ định danh nguồn gốc tài liệu ngay từ nội dung đoạn văn bản.

\paragraph{Embedding và Lưu trữ.}
Mỗi chunk được biểu diễn dưới dạng một vector 3072 chiều thông qua mô hình \texttt{text-embedding-3-large}. Nhằm tối ưu hóa thời gian xử lý corpus lớn, quá trình nhúng dữ liệu được thiết kế chạy song song đa luồng (multi-threading), batching và tích hợp cơ chế thử lại (exponential backoff) để đối phó với giới hạn rate-limit của API. Kho tri thức đầu ra được cấu trúc thành ba lớp lưu trữ đồng bộ (minh họa tại Hình~\ref{fig:pipeline}): (a) Chỉ mục vector dựa trên đồ thị HNSW~\citep{malkov2020hnsw} tối ưu cho tìm kiếm ngữ nghĩa; (b) Chỉ mục tìm kiếm toàn văn bản (full-text index) nhằm tra cứu chính xác các thực thể "cứng" như mã biểu mẫu (VD: CC01) hay số hiệu văn bản (VD: Thông tư 17/2024/TT-BCA) những trường hợp mà tìm kiếm vector thường hoạt động kém; và (c) Một catalog cơ sở dữ liệu quan hệ lưu trữ ánh xạ Tài liệu $\leftrightarrow$ Chunk $\leftrightarrow$ Vector để phục vụ việc truy vết nguồn tham chiếu.

\begin{figure}[H]
\centering
\resizebox{\textwidth}{!}{%
\begin{tikzpicture}[node distance=0.7cm and 0.9cm]
  \node[gbox] (api) {Cổng DVCQG\\(API nội bộ)};
  \node[fbox, right=of api] (crawl) {Bộ thu thập\\resume, lịch sự};
  \node[dbox, right=of crawl] (raw) {JSON thô\\5.208 thủ tục};
  \node[fbox, right=of raw] (parse) {Chuẩn hoá\\\& phân đoạn};
  \node[dbox, right=of parse] (corpus) {Corpus + metadata\\+ chunk theo mục};
  \node[fbox, right=of corpus] (embed) {Nhúng 3072 chiều\\(song song)};
  \node[abox, right=of embed] (store) {Kho tri thức\\vector + từ khoá\\+ catalog};

  \draw[farr] (api) -- (crawl);
  \draw[farr] (crawl) -- (raw);
  \draw[farr] (raw) -- (parse);
  \draw[farr] (parse) -- (corpus);
  \draw[farr] (corpus) -- (embed);
  \draw[farr] (embed) -- (store);
\end{tikzpicture}}
\caption{Pipeline thu thập dữ liệu: từ API của Cổng Dịch vụ công Quốc gia tới kho truy hồi ba lớp.}
\label{fig:pipeline}
\end{figure}

%                          
\subsection{Truy hồi và tổng hợp câu trả lời có trích dẫn}
\label{sub:rag}

Hệ thống nhận đầu vào là một câu hỏi (kèm theo lịch sử hội thoại) và trả về câu trả lời có tích hợp trích dẫn trong câu trả lời. Quá trình này vận hành qua bốn bước nối tiếp nhau (Hình~\ref{fig:rag}).

\paragraph{(1) Contextual Query Rewriting: }
Các câu hỏi từ người dùng thường ngắn, mang tính chất tiếp nối (follow-up) và thiếu ngữ cảnh (ví dụ: ``còn quy trình thì sao?'') hoặc đôi khi chỉ là lời chào giao tiếp. Ở bước này, hệ thống tận dụng lịch sử hội thoại để viết lại các câu hỏi mập mờ thành một truy vấn độc lập, trọn vẹn về mặt ngữ nghĩa (self-contained query), đồng thời chuẩn hóa các từ viết tắt phổ biến (VD: CCCD, MST, BHXH). Nếu mô hình phân loại tin nhắn chỉ là câu chào hỏi mang tính xã giao, hệ thống sẽ phát tín hiệu bỏ qua (bypass) quy trình truy hồi, cho phép mô hình ngôn ngữ phản hồi tự nhiên ngay lập tức thay vì tốn tài nguyên tra cứu.

\paragraph{(2) Hybrid Retrieval: }
Truy vấn sau khi được chuẩn hóa sẽ thực thi song song trên hai kênh: tìm kiếm ngữ nghĩa (semantic search) qua chỉ mục vector và tìm kiếm từ khoá (keyword search) qua chỉ mục toàn văn bản (full-text index). Kết quả từ hai kênh này sau đó được hợp nhất và xếp hạng lại. Sự kết hợp này khai thác triệt để ưu điểm bù trừ của hai phương pháp tiếp cận~\citep{karpukhin2020dpr,robertson2009bm25}: tìm kiếm ngữ nghĩa xử lý tốt các câu hỏi diễn đạt tự do, trong khi tìm kiếm từ khóa đảm bảo độ chính xác tuyệt đối đối với các thực thể cụ thể như mã hiệu thủ tục hay số hiệu văn bản pháp luật.

\paragraph{(3) Relevance Gating: }
Khác với chiến lược Top-$k$ truyền thống — luôn trả về $k$ tài liệu gần nhất bất kể mức độ liên quan thực tế — hệ thống của chúng tôi áp dụng một bước đánh giá trung gian. Mỗi đoạn văn bản truy hồi được một LLM đánh giá độ tương đồng ngữ nghĩa so với câu hỏi; các đoạn văn bản có điểm số dưới ngưỡng thiết lập (threshold) sẽ bị loại bỏ. Cơ chế ``cổng lọc'' (gating) này sẽ phát ra tín hiệu "thiếu thông tin" một cách tường minh nếu không có đoạn văn bản nào vượt qua bài kiểm tra độ liên quan. Tín hiệu này đóng vai trò là trigger cốt lõi để kích hoạt hành vi tự sửa lỗi (self-correction) trong kiến trúc Agentic ở Mục~\ref{sub:reasoning}.

\paragraph{(4) Generation with Citation \& Highlighting: }
Các tài liệu vượt qua cổng lọc được ghép lại làm ngữ cảnh (context) đầu vào cho LLM. Quá trình sinh văn bản được điều khiển bởi một prompt chuyên biệt, yêu cầu LLM xuất kết quả theo cấu trúc hai khối: (a) một khối liệt kê trích dẫn, gán mỗi nguồn một số thứ tự tham chiếu định dạng 【n】, đi kèm cặp tham số \texttt{START\_PHRASE}/\texttt{END\_PHRASE} trích xuất nguyên văn đoạn bằng chứng; và (b) một khối câu trả lời, trong đó số thẻ 【n】 được chèn trực tiếp vào ngay sau mỗi dữ kiện được tổng hợp. 

Dựa trên đầu ra, hệ thống đối chiếu lại cặp cụm từ trên với tài liệu gốc để highlight chính xác đoạn văn bản bằng chứng trên giao diện, đồng thời chuyển đổi các thẻ 【n】 thành liên kết ở trong câu trả lời. Nhằm kiểm soát tối đa hiện tượng ảo giác (hallucination)~\citep{Huang_2025}, prompt được thiết lập các rào cản chặt chẽ theo đặc thù lĩnh vực: chỉ sử dụng thông tin có trong ngữ cảnh, bảo toàn nguyên văn tên giấy tờ và mã biểu mẫu, đồng thời yêu cầu LLM chủ động báo cáo các phần thông tin bị thiếu thay vì tự suy diễn (extrapolation). Nếu điểm số liên quan tổng thể của ngữ cảnh ở mức thấp, hệ thống sẽ tự động hiển thị cảnh báo để người dùng chủ động kiểm chứng. Ngoài ra, một tính năng tuỳ chọn cho phép sinh sơ đồ tư duy (mind map) để tóm tắt trực quan quy trình thủ tục dưới dạng phân cấp.

\begin{figure}[H]
\centering
\resizebox{\textwidth}{!}{%
\begin{tikzpicture}[node distance=0.7cm and 0.85cm]
  \node[gbox] (q) {Câu hỏi\\+ Lịch sử};
  \node[fbox, right=of q] (rw) {(1) Contextual Query Rewriting};
  \node[fbox, right=of rw] (ret) {(2) Hybrid Retrieval};
  \node[abox, right=of ret] (gate) {(3) Relevance Gating};
  \node[fbox, right=of gate] (syn) {(4) Generation with Citation \& Highlighting};
  \node[dbox, right=of syn] (ans) {Câu trả lời 【n】\\+ Highlight + Mindmap};

  \draw[farr] (q) -- (rw);
  \draw[farr] (rw) -- (ret);
  \draw[farr] (ret) -- (gate);
  \draw[farr] (gate) -- (syn);
  \draw[farr] (syn) -- (ans);
  % nhánh phụ: chào hỏi -> trả lời thẳng; không còn đoạn liên quan -> tín hiệu thiếu
  \draw[farr, dashed] (rw.south) to[out=-90,in=-90] node[flab,below]{Chào hỏi xã giao} (ans.south);
  \draw[farr, dashed] (gate.north) to[out=90,in=90] node[flab,above]{Tín hiệu thiếu thông tin} (syn.north);
\end{tikzpicture}}
\caption{Lõi truy hồi và sinh văn bản có trích dẫn (suy luận một lượt). Đường nét đứt thể hiện luồng xử lý ngoại lệ: bỏ qua truy hồi đối với các câu chào hỏi, hoặc trả về tín hiệu thiếu thông tin khi cổng lọc loại bỏ toàn bộ tài liệu.}
\label{fig:rag}
\end{figure}

%                          
\subsection{Tích hợp các chiến lược suy luận (Reasoning Strategies)}
\label{sub:reasoning}

Kiến trúc suy luận một lượt ở Mục~\ref{sub:rag} đáp ứng tốt các truy vấn trực diện, tuy nhiên bộc lộ hai hạn chế chính khi: (i) truy vấn đòi hỏi suy luận nhiều bước (multi-step) hoặc liên quan đến nhiều thủ tục khác nhau; và (ii) thông tin người dùng yêu cầu không tồn tại trong corpus (ví dụ: địa chỉ cụ thể, giờ làm việc hay đường dây nóng của cơ quan tiếp nhận hồ sơ tại một địa phương). 

Để mở rộng năng lực hệ thống mà không phá vỡ lõi RAG (RAG core) hiện tại, chúng tôi thiết kế một interface suy luận chung. Một chiến lược (strategy) sẽ nhận đầu vào gồm (câu hỏi, lịch sử hội thoại, tập công cụ truy hồi) và xuất ra một luồng (stream) câu trả lời. Mọi chiến lược đều được implement dựa trên interface này và dùng chung module tổng hợp có trích dẫn (Mục~\ref{sub:rag}) cho khâu sinh kết quả; việc chuyển đổi giữa các chiến lược chỉ là thiết lập cấu hình runtime. Hệ thống hiện tại tích hợp ba chiến lược sau:

\paragraph{Chiến lược A: Suy luận một lượt (Simple).}
Chiến lược này tái sử dụng hoàn toàn pipeline ở Mục~\ref{sub:rag}: một lần truy hồi (retrieval) và một lần sinh (generation). Đây được xem là baseline của hệ thống với chi phí tính toán và độ trễ (latency) thấp nhất, phù hợp tối ưu cho các câu hỏi tra cứu trực diện về một thủ tục đơn lẻ.

\paragraph{Chiến lược B: Agent ReAct hai giai đoạn (Mặc định).}
Chiến lược này hiện thực hóa mô hình ReAct~\citep{yao2023react} và vận hành qua hai giai đoạn (Hình~\ref{fig:react}). Trong bối cảnh hội thoại đa lượt, truy vấn đầu vào sẽ được contextualize (ví dụ: giải tham chiếu các đại từ chỉ định như ``thủ tục này'', ``cơ quan ấy'') trước khi bước vào Giai đoạn 1.

\textbf{Giai đoạn 1 — Gom nguồn (Source Accumulation):} Agent thực thi vòng lặp \textit{Thought $\to$ Action $\to$ Observation}. Agent được trang bị hai công cụ: \texttt{docsearch} (tra cứu corpus nội bộ, tích hợp cổng lọc Relevance Gating như Mục~\ref{sub:rag}) và \texttt{web\_search} (tìm kiếm dự phòng trên Internet). Một policy "ưu tiên corpus" được thiết lập thông qua system prompt: với mọi truy vấn, hành động (Action) đầu tiên luôn phải là \texttt{docsearch}. Agent chỉ được phép kích hoạt \texttt{web\_search} khi \texttt{docsearch} trả về kết quả rỗng (không có tài liệu liên quan), hoặc khi corpus chỉ cung cấp thông tin dẫn chiếu chung chung mà truy vấn lại yêu cầu thông tin thực địa cụ thể. Cơ chế này kế thừa tư tưởng của Corrective-RAG~\citep{yan2024crag, asai2024selfrag}, được hiện thực hóa thông qua prompting kết hợp với tín hiệu từ cổng lọc thay vì phải fine-tune lại mô hình. Các đoạn văn bản từ corpus và kết quả web (được định dạng thành nguồn có thể trích dẫn và gán nhãn "Web") sẽ được đẩy vào một "bể gom nguồn" (Source Sink) dùng chung.

\textbf{Giai đoạn 2 — Tổng hợp (Synthesis):} Tập hợp nguồn dữ liệu đã gom được đưa qua module tổng hợp có trích dẫn (Mục~\ref{sub:rag}) để sinh câu trả lời kèm các thẻ trích dẫn 【n】. Các trích dẫn tham chiếu đến nguồn web sẽ được tự động dán nhãn và đính kèm hyperlink để phân biệt với nguồn chính thống. Trong suốt quá trình chạy Giai đoạn 1, từng bước \textit{Thought/Action/Observation} được stream trực tiếp lên UI. Điều này giúp người dùng cuối quan sát được lộ trình suy luận và tra cứu của agent — một yêu cầu thiết yếu về tính minh bạch (transparency) đối với các hệ thống AI trong lĩnh vực hành chính công.

\paragraph{Chiến lược C: Lập kế hoạch tách rời (ReWOO).}
Như một phương án cân bằng (trade-off), hệ thống tích hợp thêm kiến trúc ReWOO~\citep{xu2023rewoo}. Thay vì đan xen suy luận và quan sát ở mỗi bước (như ReAct), agent sẽ lập toàn bộ kế hoạch (plan) gọi công cụ từ trước, sau đó mới tiến hành thực thi (execute) và tổng hợp. Cách tiếp cận này giúp giảm đáng kể số lượt gọi LLM do không phải nạp lại toàn bộ lịch sử quan sát vào context window ở mỗi vòng lặp, đặc biệt phù hợp trong các kịch bản cần tối ưu chi phí API và độ trễ. Bảng~\ref{tab:strategy} tóm tắt sự khác biệt giữa ba chiến lược.

\begin{table}[H]

\centering

\small

\renewcommand{\arraystretch}{1.25}

\begin{tabular}{@{}lcccp{3.6cm}@{}}

\toprule

Chiến lược & Số lượt truy hồi & Web dự phòng & Tự sửa & Phù hợp \\

\midrule

Simple (một lượt) & 1 & Không & Không & Câu hỏi trực diện, một thủ tục \\

ReAct (mặc định)  & Lặp ($\leq N$) & Có & Có & Câu hỏi phức hợp; cần thông tin ngoài corpus \\

ReWOO             & Theo kế hoạch & Có & Hạn chế & Tối ưu chi phí khi dùng nhiều công cụ \\

\bottomrule

\end{tabular}

\caption{So sánh ba chiến lược suy luận. Cả ba dùng chung lõi truy hồi và lõi tổng hợp có

trích dẫn (Mục~\ref{sub:rag}); chỉ khác ở cách lập kế hoạch và gọi công cụ.}

\label{tab:strategy}

\end{table}

\begin{figure}[H]
\centering
\resizebox{\textwidth}{!}{%
\begin{tikzpicture}
  %  -- Pha 1: nodes (toạ độ tường minh để không chồng lấn)  --
  \node[gbox] (q)        at (0,0)        {Truy vấn đầu vào\\(+ Contextualize)};
  \node[fbox] (thought) at (3.3,0)       {Thought\\(Suy luận)};
  \node[fbox] (action)  at (6.6,0)       {Action\\(Chọn công cụ)};
  \node[dbox] (doc)     at (9.9,1.35)  {docsearch\\(Corpus + Gating)};
  \node[abox] (web)     at (9.9,-1.35) {web\_search\\(Dự phòng)};
  \node[fbox] (obs)     at (13.2,0)    {Observation};
  \node[gbox] (sink)    at (16.5,0)    {Source Sink\\(Corpus + Web)};
  %  -- Pha 2: nodes (hàng dưới, tách hẳn khỏi Pha 1)  --
  \node[fbox] (syn)     at (16.5,-3.7) {Giai đoạn 2: Tổng hợp\\có trích dẫn (Mục~3.2)};
  \node[dbox] (ans)     at (12.3,-3.7) {Câu trả lời 【n】\\(Gắn nhãn nguồn Web)};

  %  -- arrows: Pha 1  --
  \draw[farr] (q) -- (thought);
  \draw[farr] (thought) -- (action);
  \draw[farr] (action) -- (doc);
  \draw[farr] (action) -- (web);
  \draw[farr] (doc) -- (obs);
  \draw[farr] (web) -- (obs);
  \draw[farr] (obs) -- (sink);
  % vòng lặp: định tuyến vuông góc qua đỉnh (tránh đè lên docsearch)
  \draw[farr] (obs.north) -- ++(0,1.85) -| (thought.north);
  \node[flab] at (4.9,2.5) {Loop $\leq N$};
  %  -- arrows: Pha 1 -> Pha 2  --
  \draw[farr] (sink) -- (syn);
  \draw[farr] (syn) -- (ans);

  %  -- khung bao Pha 1 + nhãn (đặt phía trên, giữa, tách khỏi vòng lặp)  --
  \begin{scope}[on background layer]
    \node[draw=black!35, dashed, rounded corners, fill=black!3,
          fit=(thought)(action)(doc)(web)(obs), inner sep=9pt] (p1) {};
  \end{scope}
  \node[font=\footnotesize] at (8.25,3.15) {GIAI ĐOẠN 1 — Source Accumulation (Vòng lặp ReAct)};
\end{tikzpicture}}
\caption{Chiến lược Agent ReAct hai giai đoạn. Giai đoạn 1 lặp chu trình Thought/Action/Observation để gom nguồn từ corpus (ưu tiên) và web (dự phòng) theo tư tưởng Corrective-RAG; Giai đoạn 2 tái sử dụng lõi tổng hợp có trích dẫn đã định nghĩa ở Mục~\ref{sub:rag}.}
\label{fig:react}
\end{figure}

Nguyên tắc thiết kế phân tách giữa một lõi truy hồi/sinh dùng chung và các chiến lược suy luận khả cắm (pluggable reasoning strategies) mang lại hai lợi điểm lớn về mặt kiến trúc. Thứ nhất, chất lượng trích dẫn và tính trung thực (faithfulness) của câu trả lời luôn được đảm bảo nhất quán bất kể hệ thống đang chạy ở chế độ nào. Thứ hai, khả năng mở rộng (extensibility) được tối ưu: việc tích hợp thêm một framework agent mới trong tương lai chỉ yêu cầu hiện thực hóa interface chung mà không gây ảnh hưởng (side-effects) đến tầng dữ liệu hay logic sinh trích dẫn.
% ===========================================================================
% ===========================================================================
\section{Kế hoạch thực hiện}
\label{sec:plan}
% ===========================================================================

Đồ án được triển khai trong khoảng thời gian từ 29/05 đến 30/06/2026, chia thành bốn tuần với mục tiêu phát triển tiệm tiến từ khâu chuẩn bị dữ liệu \(\rightarrow\) xây dựng baseline \(\rightarrow\) phát triển Agentic RAG \(\rightarrow\) mở rộng mô hình \& hoàn thiện báo cáo. Bảng~\ref{tab:roadmap} tóm tắt lộ trình thực hiện. Trong các danh sách chi tiết bên dưới, ký hiệu \(\boxtimes\) đại diện cho hạng mục đã hoàn thành (tính đến thời điểm viết báo cáo 12/06/2026), và \(\square\) chỉ các hạng mục đang hoặc sắp thực hiện.

\begin{table}[H]
\centering
\small
\renewcommand{\arraystretch}{1.65}
\setlength{\tabcolsep}{4pt}
\begin{tabular}{@{}p{5.7cm}P{2.25cm}P{2.25cm}P{2.25cm}P{2.25cm}@{}}
\toprule
Hạng mục công việc & Tuần 1 & Tuần 2 & Tuần 3 & Tuần 4 \\
 & {\scriptsize 29/05--05/06} & {\scriptsize 06/06--12/06} & {\scriptsize 13/06--22/06} & {\scriptsize 23/06--30/06} \\
\midrule

\rowcolor{blue!12}
GĐ 1 — Dữ liệu \& Pipeline nhúng \newline
{\scriptsize Crawl, chuẩn hoá, nhúng, vector DB} &
\cellcolor{blue!55} & \cellcolor{blue!28} & \cellcolor{white} & \cellcolor{white} \\
\midrule

\rowcolor{teal!12}
GĐ 2 — Khảo sát RAG \& Baseline \newline
{\scriptsize Survey RAG, prompt domain, kiến trúc Simple} &
\cellcolor{white} & \cellcolor{teal!50} & \cellcolor{white} & \cellcolor{white} \\
\midrule

\rowcolor{orange!14}
GĐ 3 — Agentic RAG \& Tích hợp \newline
{\scriptsize ReAct 2 pha, API+UI, citation, gợi ý hội thoại} &
\cellcolor{white} & \cellcolor{orange!28} & \cellcolor{orange!62} & \cellcolor{orange!30} \\
\midrule

\rowcolor{yellow!20}
GĐ 4 — Đa LLM \& Báo cáo kỹ thuật \newline
{\scriptsize Tích hợp Azure/OpenAI/vLLM, viết PDF} &
\cellcolor{white} & \cellcolor{white} & \cellcolor{white} & \cellcolor{yellow!65} \\

\bottomrule
\end{tabular}
\caption{Lộ trình thực hiện theo tuần (29/05 -- 30/06/2026). Ô đậm = tuần trọng tâm của giai đoạn; ô nhạt = phần việc gối đầu sang tuần lân cận.}
\label{tab:roadmap}
\end{table}

\subsection{Tuần 1 (29/05 -- 05/06): Chuẩn bị Dữ liệu \& Pipeline Nhúng End-to-end}
\begin{itemize}
  \cbDone Khảo sát nguồn dữ liệu và dịch ngược (reverse-engineer) các API ẩn của Cổng Dịch vụ công Quốc gia (bao gồm endpoint danh sách và chi tiết theo công dân).
  \cbDone Xây dựng crawler tích hợp khả năng resume và cơ chế an toàn tải (rate-limiting với delay + jitter, retry kèm exponential backoff), xuất file theo cơ chế nguyên tử (atomic write); thu thập thành công 5.208 thủ tục (\texttt{data/raw/*.json}).
  \cbDone Lập trình parser chuyển đổi JSON \(\rightarrow\) Markdown: bóc tách metadata vào khối frontmatter (mã, lĩnh vực, cấp thực hiện, cơ quan, đối tượng) kèm sáu hạng mục nội dung làm sạch; xuất dữ liệu ra định dạng \texttt{chunks.jsonl} phân tách theo từng mục.
  \cbDone Cấu hình mô hình nhúng (embedding model) Azure \texttt{text-embedding-3-large} (3072 chiều), tối ưu cho ngôn ngữ tiếng Việt.
  \cbDone Nạp (ingest) corpus vào cơ sở dữ liệu lai (hybrid store) bao gồm vector DB (Chroma) và document DB (LanceDB); quá trình nạp được thực thi song song đa luồng (multi-worker) nhằm tối ưu hiệu năng.
  \cbDone Đóng gói và chia sẻ bộ chỉ mục (index) lên nền tảng HuggingFace để đảm bảo tính tái tạo (reproducibility) cho các môi trường khác mà không tiêu tốn tài nguyên nhúng lại.
\end{itemize}

\subsection{Tuần 2 (06/06 -- 12/06): Khảo sát RAG \& Triển khai Kiến trúc Baseline}
\begin{itemize}
  \cbDone Khảo sát toàn diện các kiến trúc RAG (Naive/Advanced/Modular; Retrieve-Rerank; Query Rewriting; Hybrid Search) và lựa chọn thiết kế tối ưu với đặc thù miền dữ liệu.
  \cbDone Thiết kế hệ thống System Prompt QA tiếng Việt chuyên biệt (đảm bảo tính toàn vẹn của tên giấy tờ/mã mẫu/số văn bản; yêu cầu liệt kê rõ thành phần hồ sơ kèm số lượng bản chính/bản sao).
  \cbDone Tích hợp module Query Rewriting: giải nghĩa các từ viết tắt phổ biến (CCCD, MST, BHXH\ldots) và thiết lập cơ chế rẽ nhánh (bypass) bỏ qua các câu chào hỏi xã giao.
  \cbDone Triển khai kiến trúc Baseline Simple tích hợp lên UI theo pipeline: Query Rewriting \(\rightarrow\) Hybrid Retrieval \(\rightarrow\) Reranking \(\rightarrow\) Sinh câu trả lời kèm nguồn.
  \cbTodo Đánh giá định tính (qualitative evaluation) kiến trúc baseline trên một tập truy vấn mẫu để thiết lập mốc cơ sở (benchmark).
\end{itemize}

\subsection{Tuần 3 (13/06 -- 22/06): Agentic RAG, Tích hợp \& Tối ưu Tính Minh bạch}
\begin{itemize}
  \cbDone Thiết kế kiến trúc Agentic RAG hai giai đoạn dựa trên framework ReAct.
  \cbDone Giai đoạn 1: Triển khai vòng lặp ReAct với công cụ \texttt{docsearch} (tra cứu corpus, độ ưu tiên cao, tích hợp cổng lọc Relevance Gating) và \texttt{web\_search} (tra cứu dự phòng qua Brave Search, hỗ trợ trích xuất nội dung trang).
  \cbDone Tối ưu Prompt ReAct tiếng Việt: giữ nguyên các từ khóa cấu trúc Thought/Action/Observation; thiết lập policy "ưu tiên corpus" và chỉ gọi web khi corpus cung cấp thông tin dẫn chiếu thiếu tính cụ thể.
  \cbDone Giai đoạn 2: Tổng hợp câu trả lời với trích dẫn nội dòng (inline citation) định dạng \texttt{【n】}, đính kèm tham số START/END\_PHRASE để phục vụ cơ chế tô sáng (highlight) đoạn bằng chứng gốc.
  \cbDone Phân loại và dán nhãn cảnh báo đối với các nguồn web chưa qua thẩm định ngay trên panel giao diện.
  \cbDone Trực quan hoá lộ trình suy luận: truyền phát (streaming) các bước Thought/Action của Agent tới UI theo thời gian thực thông qua kết nối SSE.
  \cbDone Tích hợp tính năng gợi ý câu hỏi tiếp nối (follow-up questions) dựa trên bối cảnh hội thoại.
  \cbDone Hỗ trợ xử lý ngữ cảnh đa lượt: Contextualize câu hỏi (giải tham chiếu các đại từ chỉ định như "thủ tục này", "nó"\ldots).
  \cbDone Tự động sinh sơ đồ tư duy (mind map) phân cấp thủ tục bằng PlantUML, gom nhóm trực quan các loại giấy tờ bắt buộc/tuỳ chọn.
  \cbDone Tích hợp toàn diện backend FastAPI (hỗ trợ streaming) và frontend React (trang bị khung chat, panel nguồn, bảng log suy luận, ô gợi ý, và điều khiển regen/stop).
  \cbTodo Đánh giá hệ thống Agentic đối chiếu với Baseline (về độ chính xác, độ phủ nguồn, tính trung thực) và tiến hành phân tích lỗi (error analysis).
\end{itemize}

\subsection{Tuần 4 (23/06 -- 30/06): Tích hợp Đa LLM (Multi-LLM) \& Hoàn thiện Báo cáo}
\begin{itemize}
  \cbTodo Trừu tượng hoá (Abstract) tầng giao tiếp LLM và Embedding thông qua file cấu hình, loại bỏ sự phụ thuộc cứng (hard-code) vào một nhà cung cấp duy nhất.
  \cbTodo Xây dựng các adapter giao tiếp cho Azure OpenAI (hiện hành), OpenAI chuẩn, và các mô hình mã nguồn mở được phục vụ qua vLLM (hỗ trợ API tương thích OpenAI).
  \cbTodo Cho phép linh hoạt chuyển đổi mô hình (model switching) thông qua biến môi trường \texttt{.env} hoặc UI Settings mà không cần can thiệp mã nguồn.
  \cbTodo Kiểm thử tích hợp End-to-end với từng nhà cung cấp dịch vụ (service provider).
  \cbTodo Hoàn thiện báo cáo kỹ thuật PDF (tài liệu hoá kiến trúc, pipeline, thiết lập môi trường, hướng dẫn tái tạo) và kiện toàn phần Methodology của báo cáo hiện tại.
  \cbTodo Báo cáo đồ án môn học. 
\end{itemize}

\clearpage
\bibliographystyle{plainnat}
\bibliography{Mendeley}
\end{document}
